\section{The Early Days: Neurons, Winters, and a Long Wait}

\subsection{The First Spark (1950s--1960s)}

The story begins with a deceptively simple question: \emph{Can we build a machine that learns the way the brain does?}

In 1943, Warren McCulloch and Walter Pitts proposed the first mathematical model of a neuron --- a simple unit that takes inputs, applies a threshold, and fires an output \cite{mcculloch1943logical}.
In 1958, Frank Rosenblatt built on this idea to create the \textbf{Perceptron}: a single-layer model that could learn to classify inputs by adjusting numerical weights \cite{rosenblatt1958perceptron}.

The excitement was immediate.
The New York Times reported that the Navy had built a machine that could ``learn, make decisions and translate languages.''
It felt like the dawn of thinking machines.

\begin{tcolorbox}[colback=blue!5!white, colframe=blue!40!black, title=What is a Perceptron?]
A perceptron takes a set of numbers as input (say, pixel values of an image), multiplies each by a \emph{weight} (its importance), adds them up, and checks whether the total crosses a threshold.
If it does, the perceptron fires a ``yes''; otherwise, a ``no.''
Training adjusts the weights based on mistakes, nudging the model toward correct answers over time.
\end{tcolorbox}

\subsection{The First AI Winter (1970s)}

The optimism did not last.
In 1969, Marvin Minsky and Seymour Papert published \emph{Perceptrons} \cite{minsky1969perceptrons}, a mathematical analysis that showed single-layer perceptrons had fundamental limitations.
They could not learn even simple non-linear functions like XOR.

Funding dried up.
Progress stalled.
The field entered what is now called the \textbf{first AI winter} --- a period of disillusionment and neglect that lasted through much of the 1970s.

\subsection{The Backpropagation Revival (1980s)}

The thaw came in 1986.
David Rumelhart, Geoffrey Hinton, and Ronald Williams published a landmark paper showing how to train \emph{multi-layer} networks using an algorithm called \textbf{backpropagation} \cite{rumelhart1986learning}.

The key insight was elegant: if you can measure how wrong your model's output is, you can trace that error backwards through every layer of the network and adjust each weight accordingly.
Layer by layer, from output back to input, the error signal propagates and the network learns.

\begin{tcolorbox}[colback=green!5!white, colframe=green!40!black, title=Backpropagation in Plain English]
Imagine you are trying to throw a ball into a basket from far away.
You throw, miss, and adjust --- maybe throw a bit higher, a bit more to the left.
Backpropagation is the mathematical version of that correction signal, applied across every connection in the network simultaneously.
\end{tcolorbox}

Multi-layer networks could now learn non-linear patterns.
The XOR problem was no longer a barrier.
But the field still faced practical limits: computers were too slow, datasets too small, and training deep networks took weeks even for modest tasks.

\subsection{The Second AI Winter (1990s)}

A second winter followed.
Support Vector Machines and other classical techniques often outperformed neural networks on benchmark tasks, and they were far cheaper to train.
Neural networks fell out of fashion again.
Many researchers moved on.

A small group --- Geoffrey Hinton, Yann LeCun, Yoshua Bengio, and a few others --- kept working.
They would eventually be called the \textbf{``Godfathers of Deep Learning''} and share the Turing Award in 2018.
But in the 1990s, they were largely swimming against the current.
